\SourceLeaf{89}%

Nehmen wir nun an, die Funktionen $fx$ und $\varphi x$ hängen so miteinander zusammen, dass

\[ \Phi x = \psi x \mathord{.} Fx \]

\noindent oder, was dasselbe ist, dass

\[ \varphi x = \int \psi x \mathord{.} \frac{dfx}{dx}\,dx, \]

\noindent und berücksichtigen wir, dass

\[ \frac{dfx}{du} = \psi x\,\frac{dfx}{dt}, \]

\noindent so ergibt sich

\[ \frac{dfx}{du} = \frac{d\varphi x}{dt}. \]

\noindent Mithin [folgt], da

\[ \frac{d\varphi t}{dt} = \Phi t = \psi t \mathord{.} Ft, \]

\[\begin{aligned}
T_1 + 2 T_2 u + 3 T_3 uu + \text{u.\,s.\,w.}
&= \psi t\mathord{.}Ft+\frac{d\Theta_1}{dt}u\\
&\quad+\frac{d\Theta_2}{dt}uu+\frac{d\Theta_3}{dt}u^3+\text{u.\,s.\,w.}
\end{aligned}\]

\noindent Demnach erhalten wir zunächst

\[ T_1 = \psi t \mathord{.} Ft. \]

\noindent Da nach der obigen Bemerkung $T_1$ in $\Theta_1$ übergehen muss, wenn man im Ausdruck für die erstere Funktion an die Stelle von $Ft$ nun $\Phi t$ setzt, d.\,h. $\psi t \mathord{.} Ft$, ergibt sich

\[ \Theta_1 = (\psi t)^2 \mathord{.} Ft. \]

\noindent Daraus ergibt sich weiter

\[ T_2 = \frac{1}{2}\,\frac{d(\psi t)^2 Ft}{dt}, \]

\noindent und dann wieder nach der obigen Bemerkung

\[ \Theta_2 = \frac{1}{2}\,\frac{d(\psi t)^3 Ft}{dt}. \]

\noindent Sodann

\[ T_3 = \frac{1}{2.3}\,\frac{dd(\psi t)^3 Ft}{dt^2}, \]
\[ \Theta_3 = \frac{1}{2.3}\,\frac{dd(\psi t)^4 Ft}{dt^2}, \]
\[ T_4 = \frac{1}{2.3.4}\,\frac{d^3(\psi t)^4 Ft}{dt^3}, \]

\begin{center}usw.\end{center}
